\documentclass[11pt]{article}
\usepackage[utf8]{inputenc}
\usepackage[T1]{fontenc}
\usepackage[margin=1in]{geometry}
\usepackage{amsmath,amssymb,amsthm}
\usepackage{booktabs}
\usepackage{hyperref}
\usepackage[numbers]{natbib}
\usepackage{xcolor}

\newtheorem{proposition}{Proposition}
\newtheorem{lemma}{Lemma}

\title{Supplementary Material for Review\\
\large Geodesic Diagnostic Trajectories (GDT):\\
A Differential Geometry Framework for Adaptive Decision Support\\
and Operational Optimization in EHR Systems}
\author{Hemanth Manchabale Papachappa \and Aniriuddha Ganguly}
\date{September 2026}

\begin{document}
\maketitle
\tableofcontents
\newpage

% ─────────────────────────────────────────────────────────────
\section{Extended Proof of Proposition~1}
\label{sec:proof_extended}

The proof in the main manuscript establishes the bound via the triangle inequality
in a single step. We provide here a more detailed derivation that clarifies the
role of each term and examines the tightness of the bound.

\subsection{Setup and Notation}
Let $\xi \in \mathbb{R}^d$ with $\|\xi\| = 1$ be the unit vector spanning the
stale-topic subspace. The orthogonal projector is $P_\xi = \xi\xi^\top$. The
gated embedding is:
\begin{equation}
  \tilde{e}_t = (1 - g_t)\,h_{t-1} + g_t\,e_t,
  \quad g_t = \sigma\!\left(\frac{\kappa_t - \kappa_0}{\tau}\right) \in (0,1).
  \label{eq:gate_supp}
\end{equation}

\subsection{Formal Statement}
\begin{proposition}[Context Interference Bound --- Extended]
\label{prop:interference_extended}
Under the parameterization in Eq.~\eqref{eq:gate_supp}:
\begin{equation}
  \|P_\xi\,\tilde{e}_t\| \leq (1-g_t)\,\|P_\xi h_{t-1}\| + g_t\,\|P_\xi e_t\|.
  \label{eq:bound_supp}
\end{equation}
Equality holds if and only if $P_\xi h_{t-1}$ and $P_\xi e_t$ are co-aligned,
i.e., $P_\xi h_{t-1} = \lambda P_\xi e_t$ for some $\lambda \geq 0$.
\end{proposition}

\begin{proof}
By the linearity of $P_\xi$ and the triangle inequality on $\mathbb{R}^d$:
\begin{align}
  \|P_\xi\,\tilde{e}_t\|
  &= \|P_\xi\bigl[(1-g_t)h_{t-1} + g_t e_t\bigr]\| \nonumber \\
  &= \|(1-g_t)P_\xi h_{t-1} + g_t P_\xi e_t\| \nonumber \\
  &\leq (1-g_t)\|P_\xi h_{t-1}\| + g_t\|P_\xi e_t\|.
\end{align}
The equality condition for the triangle inequality requires $P_\xi h_{t-1}$ and
$P_\xi e_t$ to point in the same direction, i.e., $P_\xi h_{t-1} = \lambda P_\xi e_t$
with $\lambda \geq 0$.
\end{proof}

\subsection{Tightness Analysis}
When $g_t \to 1$ (full pivot suppression), the bound reduces to $\|P_\xi e_t\|$,
the projection of the \emph{current} query's embedding onto the stale-topic
subspace. This quantity is empirically small when the query encoder successfully
represents the new topic. In our benchmark, the mean value of $\|P_\xi e_t\|$ at
gate-triggered turns was $0.042 \pm 0.019$ normalized units, confirming that the
bound is tight and the post-gate interference is practically negligible.

When $g_t \to 0$ (no pivot, history trusted), the bound reduces to
$\|P_\xi h_{t-1}\|$, which equals the interference from the unmodified history
embedding---consistent with the gate being ``closed.''

\subsection{Empirical Verification Detail}
For each of the 1,074 gate-triggered turns across 31 vignettes, we computed:
\begin{itemize}
  \item \textbf{Actual interference:} $A_t = \|P_\xi \tilde{e}_t\|$, where
        $\xi$ was taken as the top eigenvector of the prior turn's SPD matrix
        $S_{t-1}$ (the principal direction of the stale intent).
  \item \textbf{Bound value:} $B_t = (1-g_t)\|P_\xi h_{t-1}\| + g_t\|P_\xi e_t\|$.
  \item \textbf{Margin:} $M_t = B_t - A_t \geq 0$ (positive in all 1,074 cases).
\end{itemize}
Summary statistics: $\min(M_t) = 0.003$, $\max(M_t) = 0.089$, $\bar{M} = 0.031$,
$\sigma(M) = 0.018$ normalized units. The bound was never violated.

% ─────────────────────────────────────────────────────────────
\section{Hyperparameter Settings}
\label{sec:hyperparams}

\begin{table}[h]
\centering
\caption{GDT hyperparameters used in all reported experiments.}
\label{tab:hyperparams}
\begin{tabular}{llll}
\toprule
\textbf{Parameter} & \textbf{Symbol} & \textbf{Value} & \textbf{Notes} \\
\midrule
Embedding dimension    & $d$       & 128    & Post-projection (ClinicalBERT: 768-d) \\
Sliding window         & $k$       & 5      & Turns for covariance computation \\
Regularization         & $\lambda$ & $10^{-4}$ & Ensures SPD condition \\
Gate threshold         & $\kappa_0$& 1.2    & Pivot detection threshold \\
Gate temperature       & $\tau$    & 0.3    & Sigmoid sharpness \\
Confidence scale       & $\alpha$  & 2.0    & Prefetch decay after pivot \\
JEPA predictor         & ---       & 3-layer MLP & Hidden dim 256, ReLU \\
Momentum coefficient   & $m$       & 0.996  & Exponential moving average \\
Prefetch top-$k$       & ---       & 10     & Documents retrieved per turn \\
Learning rate          & ---       & $3\times10^{-4}$ & AdamW, cosine schedule \\
Batch size             & ---       & 32     & Sessions per batch \\
Training epochs        & ---       & 50     & Early stopping (patience 5) \\
\bottomrule
\end{tabular}
\end{table}

% ─────────────────────────────────────────────────────────────
\section{Dataset Details}
\label{sec:datasets}

\begin{table}[h]
\centering
\caption{Clinical datasets used in vignette construction.}
\label{tab:datasets}
\begin{tabular}{lllll}
\toprule
\textbf{Dataset} & \textbf{Version} & \textbf{Size} & \textbf{Specialty mix} & \textbf{Access} \\
\midrule
MIMIC-III & v1.4  & 46,520 admissions & ICU (mixed) & PhysioNet \\
MIMIC-IV  & v2.2  & 65,366 admissions & ICU + ED    & PhysioNet \\
eICU      & v2.0  & 200,859 admissions & ICU        & PhysioNet \\
HiRID     & v1.1.1& 33,905 admissions & ICU        & PhysioNet \\
\bottomrule
\end{tabular}
\end{table}

\noindent The 31 test vignettes were distributed as follows across specialty areas:
critical care ($n{=}11$, 35\%), emergency medicine ($n{=}8$, 26\%), hospital
medicine ($n{=}7$, 23\%), internal medicine ($n{=}5$, 16\%). Each vignette required
a minimum of 3 abrupt topic pivots (median 4, range 3--7).

% ─────────────────────────────────────────────────────────────
\section{Full Statistical Results}
\label{sec:stats_full}

\begin{table}[h]
\centering
\caption{Complete pairwise comparison results (Wilcoxon signed-rank, two-tailed).}
\label{tab:stats}
\begin{tabular}{llrrl}
\toprule
\textbf{Contrast} & \textbf{Outcome} & \textbf{$p$-value} & \textbf{Cohen's $d$} & \textbf{Interpretation} \\
\midrule
GDT vs.\ TF-IDF & Latency     & $<0.001$ & $-4.12$ & Large effect \\
GDT vs.\ BM25   & Latency     & $<0.001$ & $-4.10$ & Large effect \\
GDT vs.\ CMA    & Latency     & $<0.001$ & $-0.87$ & Moderate effect \\
CMA vs.\ TF-IDF & Latency     & $<0.001$ & $-3.91$ & Large effect \\
BM25 vs.\ TF-IDF& Latency     & $0.68$   & $+0.01$ & Negligible \\
\midrule
GDT vs.\ TF-IDF & NASA-TLX    & $0.003$  & $-0.55$ & Medium effect \\
GDT vs.\ BM25   & NASA-TLX    & $0.007$  & $-0.52$ & Medium effect \\
GDT vs.\ CMA    & NASA-TLX    & $0.041$  & $-0.19$ & Small effect \\
CMA vs.\ TF-IDF & NASA-TLX    & $0.007$  & $-0.49$ & Medium effect \\
\midrule
All contrasts    & Time        & $\geq 0.35$ & $\leq 0.03$ & Non-significant \\
All contrasts    & Task compl. & N/A      & N/A     & 0\% all arms \\
\bottomrule
\end{tabular}
\end{table}

% ─────────────────────────────────────────────────────────────
\section{M/M/c Queueing Model --- Sensitivity Analysis}
\label{sec:queueing_supp}

The main text reports a 4.1-minute per-reviewer reduction under one specific
parameter setting ($c{=}10$, $\lambda_q{=}8$ queries/min). Table~\ref{tab:queueing_sensitivity}
provides the full sensitivity analysis.

\begin{table}[h]
\centering
\caption{Projected per-reviewer time savings (minutes per 4-hour shift) under
varying M/M/c parameters. GDT vs.\ TF-IDF; $\mu_q$ computed from observed
mean latencies.}
\label{tab:queueing_sensitivity}
\begin{tabular}{rrrrr}
\toprule
\textbf{Reviewers ($c$)} & \textbf{$\lambda_q$ (q/min)} & \textbf{$\rho/c$} & \textbf{$\Delta W$ (ms)} & \textbf{Savings (min/shift)} \\
\midrule
5  & 4  & 0.52 & $-47.1$ & 1.9 \\
5  & 8  & 0.61 & $-49.8$ & 2.4 \\
10 & 8  & 0.48 & $-50.8$ & 4.1 \\
10 & 16 & 0.59 & $-52.3$ & 5.0 \\
20 & 16 & 0.47 & $-51.0$ & 4.9 \\
\bottomrule
\end{tabular}
\end{table}

The projected savings are robust across the parameter range, varying between 1.9
and 5.0 minutes per reviewer per shift. The offered load $\rho/c < 1$ in all
reported scenarios, confirming system stability under the M/M/c assumptions.

% ─────────────────────────────────────────────────────────────
\section{Ablation: Gate Threshold Sensitivity}
\label{sec:gate_sensitivity}

To assess sensitivity to the gate threshold $\kappa_0$, we swept $\kappa_0 \in
\{1.0, 1.2, 1.5, 1.8\}$ on the held-out vignettes while holding all other
hyperparameters fixed.

\begin{table}[h]
\centering
\caption{Latency and trigger rate as a function of gate threshold $\kappa_0$.}
\label{tab:gate_sensitivity}
\begin{tabular}{rrrr}
\toprule
$\kappa_0$ & \textbf{Mean Latency (ms)} & \textbf{Trigger Rate (\%)} & \textbf{$\Delta$ vs.\ TF-IDF} \\
\midrule
1.0 & 136.1 & 41.2 & $-$28.0\% \\
1.2 & 138.2 & 34.7 & $-$26.9\% \\
1.5 & 140.8 & 26.3 & $-$25.5\% \\
1.8 & 143.1 & 18.9 & $-$24.3\% \\
\bottomrule
\end{tabular}
\end{table}

Results are stable across the range. $\kappa_0{=}1.2$ was selected as the
default because it balances trigger sensitivity (34.7\% of transitions, consistent
with our estimate of multi-pivot frequency in clinical sessions) against latency.
Lower thresholds trigger more aggressively and reduce latency slightly but may
suppress context during within-topic refinements.

\bibliographystyle{unsrtnat}
\bibliography{../references}

\end{document}
